% GENERATED FILE. Edit canonical lessons, then run npm run generate:books.
% canonical-source-hash: fnv1a64:5298181e6f0fde80
% canonical-lessons: GE-C364-bedingung, GE-C364-voraussetzung, GE-C364-voraussetzen, GE-C364-sofern, GE-C364-soweit

\chapter{Unter einer Bedingung}
\label{ch:ge-unter-einer-bedingung}
\hlchaptermodality{364}

\begin{chapteropening}
\textbf{By the end of this chapter:} \emph{I can set a condition, name what has to be true first, and say provided that with sofern and as far as with soweit.}

Complete the last lesson of chapter 364: I can set a condition, name what has to be true first, and say provided that with sofern and as far as with soweit.
\end{chapteropening}

\section[die Bedingung]{die Bedingung — the condition}

\label{lesson:GE-C364-bedingung}



\begin{quote}
Before the new one: say the German for to recycle, then the German for to consume.
\end{quote}

\subsection*{You'll want to know: die Bedingung}
\textbf{die Bedingung} — \textquotedblleft{}the condition\textquotedblright{}.

\begin{quote}
Ich komme mit, unter einer Bedingung: Du bezahlst das Taxi.
\end{quote}

\begin{quote}
\emph{I will come along on one condition: you pay for the taxi.}
\end{quote}

Plural: \emph{die Bedingungen}. \emph{Unter der Bedingung, dass …} — on condition that …

\begin{cousinweb}
\emph{Bedingung} is the noun of \emph{bedingen} \textquotedblleft{}to stipulate\textquotedblright{}, inherited from Middle High German \emph{bedingen} \textquotedblleft{}to hire, to engage\textquotedblright{} (EtymDB). The verb is usually explained as built on \emph{Ding} in its old sense of a court or assembly, Old High German \emph{ding}, Proto-Germanic \emph{þingą} (EtymDB): terms agreed before the assembly.
\end{cousinweb}

\begin{grammarlens}[title={What you've built}]
One more word for telling, explaining and describing at B1.
\end{grammarlens}

\subsection*{Your turn}
\begin{itemize}
\raggedright
  \item \emph{Say it:} \emph{die Bedingung}
  \item \emph{Say it:} \emph{die Bedingung}, once more
  \item \emph{Say it:} a sentence of your own with \emph{die Bedingung}
  \item \emph{Recall:} say the German for to recycle, then the German for to consume, then say \emph{die Bedingung} again
\end{itemize}

\begin{tcolorbox}[breakable,colback=teal!4,colframe=teal!35!black,title={Before you move on}]
What does die Bedingung mean? (\textquotedblleft{}The condition\textquotedblright{}.) Say it once more.
\end{tcolorbox}
\section[die Voraussetzung]{die Voraussetzung — the precondition, the requirement}

\label{lesson:GE-C364-voraussetzung}



\begin{quote}
Before the new one: say the German for to consume, then the German for the condition.
\end{quote}

\subsection*{You'll want to know: die Voraussetzung}
\textbf{die Voraussetzung} — \textquotedblleft{}the precondition, the requirement\textquotedblright{}.

\begin{quote}
Gute Deutschkenntnisse sind eine Voraussetzung für die Stelle.
\end{quote}

\begin{quote}
\emph{A good knowledge of German is a requirement for the job.}
\end{quote}

Plural: \emph{die Voraussetzungen}. \emph{Unter der Voraussetzung, dass …} — provided that …

\begin{cousinweb}
\emph{Voraussetzung} is the noun of \emph{voraussetzen}, built from \emph{voraus} \textquotedblleft{}ahead\textquotedblright{} and \emph{setzen} \textquotedblleft{}to set\textquotedblright{}, Proto-Germanic \emph{satjaną} (EtymDB): what is set in place beforehand.
\end{cousinweb}

\begin{grammarlens}[title={What you've built}]
One more word for telling, explaining and describing at B1.
\end{grammarlens}

\subsection*{Your turn}
\begin{itemize}
\raggedright
  \item \emph{Say it:} \emph{die Voraussetzung}
  \item \emph{Say it:} \emph{die Voraussetzung}, once more
  \item \emph{Say it:} a sentence of your own with \emph{die Voraussetzung}
  \item \emph{Recall:} say the German for to consume, then the German for the condition, then say \emph{die Voraussetzung} again
\end{itemize}

\begin{tcolorbox}[breakable,colback=teal!4,colframe=teal!35!black,title={Before you move on}]
What does die Voraussetzung mean? (\textquotedblleft{}The precondition, the requirement\textquotedblright{}.) Say it once more.
\end{tcolorbox}
\section[voraussetzen]{voraussetzen — to require, to presuppose; to take for granted}

\label{lesson:GE-C364-voraussetzen}



\begin{quote}
Before the new one: say the German for the condition, then the German for the precondition.
\end{quote}

\subsection*{You'll want to know: voraussetzen}
\textbf{voraussetzen} — \textquotedblleft{}to require, to presuppose; to take for granted\textquotedblright{}.

\begin{quote}
Die Stelle setzt keine Vorkenntnisse voraus.
\end{quote}

\begin{quote}
\emph{The job requires no previous knowledge.}
\end{quote}

Separable: \emph{das setzt … voraus}, \emph{das hat … vorausgesetzt}. \emph{Vorausgesetzt, dass …} — provided that …

\begin{cousinweb}
\emph{voraussetzen} is built from \emph{voraus} \textquotedblleft{}ahead\textquotedblright{} and \emph{setzen} \textquotedblleft{}to set\textquotedblright{}, inherited from Proto-Germanic \emph{satjaną} (EtymDB), a cousin of English \emph{set}.
\end{cousinweb}

\begin{grammarlens}[title={What you've built}]
One more word for telling, explaining and describing at B1.
\end{grammarlens}

\subsection*{Your turn}
\begin{itemize}
\raggedright
  \item \emph{Say it:} \emph{voraussetzen}
  \item \emph{Say it:} \emph{voraussetzen}, once more
  \item \emph{Say it:} a sentence of your own with \emph{voraussetzen}
  \item \emph{Recall:} say the German for the condition, then the German for the precondition, then say \emph{voraussetzen} again
\end{itemize}

\begin{tcolorbox}[breakable,colback=teal!4,colframe=teal!35!black,title={Before you move on}]
What does voraussetzen mean? (\textquotedblleft{}To require, to presuppose; to take for granted\textquotedblright{}.) Say it once more.
\end{tcolorbox}
\section[sofern]{sofern — provided that, as long as}

\label{lesson:GE-C364-sofern}



\begin{quote}
Before the new one: say the German for the precondition, then the German for to require.
\end{quote}

\subsection*{You'll want to know: sofern}
\textbf{sofern} — \textquotedblleft{}provided that, as long as\textquotedblright{}.

\begin{quote}
Wir grillen am Samstag, sofern das Wetter gut ist.
\end{quote}

\begin{quote}
\emph{We are having a barbecue on Saturday, provided the weather is good.}
\end{quote}

\emph{Sofern} sends the verb to the end, like \emph{wenn}: \emph{…, sofern das Wetter gut ist}. It sounds a little more formal.

\begin{cousinweb}
\emph{sofern} joins \emph{so} and \emph{fern} \textquotedblleft{}far\textquotedblright{}, from Old High German \emph{ferrana} (EtymDB): \textquotedblleft{}so far as\textquotedblright{}, to the extent that something holds.
\end{cousinweb}

\begin{grammarlens}[title={What you've built}]
One more word for telling, explaining and describing at B1.
\end{grammarlens}

\subsection*{Your turn}
\begin{itemize}
\raggedright
  \item \emph{Say it:} \emph{sofern}
  \item \emph{Say it:} \emph{sofern}, once more
  \item \emph{Say it:} a sentence of your own with \emph{sofern}
  \item \emph{Recall:} say the German for the precondition, then the German for to require, then say \emph{sofern} again
\end{itemize}

\begin{tcolorbox}[breakable,colback=teal!4,colframe=teal!35!black,title={Before you move on}]
What does sofern mean? (\textquotedblleft{}Provided that, as long as\textquotedblright{}.) Say it once more.
\end{tcolorbox}
\section[soweit]{soweit — as far as}

\label{lesson:GE-C364-soweit}



\begin{quote}
Before the new one: say the German for to require, then the German for provided that.
\end{quote}

\subsection*{You'll want to know: soweit}
\textbf{soweit} — \textquotedblleft{}as far as\textquotedblright{}.

\begin{quote}
Soweit ich weiß, ist das Geschäft am Montag geschlossen.
\end{quote}

\begin{quote}
\emph{As far as I know, the shop is closed on Monday.}
\end{quote}

\emph{Soweit} sends the verb to the end: \emph{soweit ich weiß}, \emph{soweit ich mich erinnere}.

\begin{cousinweb}
\emph{soweit} joins \emph{so} and \emph{weit} \textquotedblleft{}wide, far\textquotedblright{}, inherited from Middle High German \emph{wīt}, Old High German \emph{wīt} and Proto-Germanic \emph{wīdaz} (EtymDB), a cousin of English \emph{wide}.
\end{cousinweb}

\begin{grammarlens}[title={What you've built}]
One more word for telling, explaining and describing at B1.
\end{grammarlens}

\subsection*{Your turn}
\begin{itemize}
\raggedright
  \item \emph{Say it:} \emph{soweit}
  \item \emph{Say it:} \emph{soweit}, once more
  \item \emph{Say it:} a sentence of your own with \emph{soweit}
  \item \emph{Recall:} say the German for to require, then the German for provided that, then say \emph{soweit} again
\end{itemize}

\begin{tcolorbox}[breakable,colback=teal!4,colframe=teal!35!black,title={Before you move on}]
What does soweit mean? (\textquotedblleft{}As far as\textquotedblright{}.) Say it once more.
\end{tcolorbox}
